\documentclass[../../main.tex]{subfiles}% 注意这里的文件路径不能用 ./main.tex ,否则用latexmk编译子文件会报错
\graphicspath{{\subfix{./image/}}{\subfix{../../image/}}} % 指定图片引用路径,后续可以直接使用图片文件名
% 如果图片存放在上级目录,则使用 ../../image/ 作为备用路径

% 例如:
% \begin{figure}[H]
% \centering
% \includegraphics[width=0.3\textwidth]{图.png}
% \caption{}
% \label{figure:图}
% \end{figure}
% 注意:上述\label{}一定要放在\caption{}之后,否则引用图片序号会只会显示??.

\begin{document}

\section{草稿三}

\begin{theorem}[扩张定理]\label{theorem:ch05-扩张定理}
设 $1\leqslant p\leqslant\infty$, $U\subseteq\mathbb{R}^{n}$ 是有界开集, 且 $\partial U$ 为 $C^{1}$. 设 $V$ 是满足 $\overline{U}\subseteq V$ 的有界开集. 则存在有界线性算子
\begin{align*}
E:W^{1,p}(U)\longrightarrow W^{1,p}(\mathbb{R}^{n}),
\end{align*}
使得对每个 $u\in W^{1,p}(U)$, 都有:
\begin{enumerate}[(1)]
\item $Eu=u$ 在 $U$ 上几乎处处成立.
\item 存在仅由 $U,V$ 确定的紧集 $K_*\subseteq V$, 使得 $\operatorname{spt}(Eu)\subseteq K_*$.
\item 存在仅依赖于 $n,p,U,V$ 的常数 $C>0$, 使得
\begin{align*}
\|Eu\|_{W^{1,p}(\mathbb{R}^{n})}
\leqslant C\|u\|_{W^{1,p}(U)}.
\end{align*}
\end{enumerate}
这里 Sobolev 函数的支集按本性支集理解. 称 $Eu$ 为 $u$ 到 $\mathbb{R}^{n}$ 的扩张.
\end{theorem}

\begin{figure}[H]
\centering
% 在此添加原图 ch05-half-ball.png.
\caption{边界处的半球}
\label{figure:ch05-边界半球}
\end{figure}

\begin{figure}[htbp]
\centering
% 在此添加原图 8e1ffee58f4e42d8a67b7ea75f26d99d.png.
\caption{拉直边界}
\label{figure:ch05-拉直边界}
\end{figure}

\begin{proof}
若 $U=\varnothing$, 结论显然成立. 以下设 $U\neq\varnothing$.

证明中将标量函数的梯度 $Df$ 视为行向量, $|Df|$ 表示其欧氏长度, 并采用以下等价的 Sobolev 范数:
\begin{align*}
\|f\|_{W^{1,p}(A)}
&\triangleq
\left(
\|f\|_{L^{p}(A)}^{p}
+\|Df\|_{L^{p}(A)}^{p}
\right)^{\frac{1}{p}},
\quad 1\leqslant p<\infty,\\
\|f\|_{W^{1,\infty}(A)}
&\triangleq
\|f\|_{L^{\infty}(A)}
+\|Df\|_{L^{\infty}(A)}.
\end{align*}
所有边界坐标邻域及单位分解均预先固定, 不随被扩张的函数而改变.

\noindent$\bm{Step}\bm{1}$.{\heiti (构造平坦边界上的反射扩张)}

设
\begin{align*}
B&\triangleq B(0,\rho),\\
B^+&\triangleq B\cap\{x_n>0\},\\
B^-&\triangleq B\cap\{x_n<0\},\\
\Sigma&\triangleq B\cap\{x_n=0\}.
\end{align*}
这里 $B^+$ 和 $B^-$ 都是开半球, 平坦边界为 $\Sigma$, 见\cref{figure:ch05-边界半球}.

先设 $f\in C^{1}(\overline{B^+})$. 定义
\begin{align*}
R_1(x',x_n)\triangleq(x',-x_n),\quad
R_2(x',x_n)\triangleq\left(x',-\frac{x_n}{2}\right).
\end{align*}
若 $x\in B^-$, 则 $R_1x,R_2x\in B^+$. 定义反射算子
\begin{align}
(\mathcal{R}f)(x)
\triangleq
\begin{cases}
f(x),&x\in B^+\cup\Sigma,\\
-3f(R_1x)+4f(R_2x),&x\in B^-.
\end{cases}
\label{align:ch05-反射延拓}
\end{align}
由于 $-3+4=1$, 当下半球中的点趋向 $\Sigma$ 上的点时, 上式的第二部分趋向对应的 $f$ 值, 因而 $\mathcal{R}f$ 在整个 $B$ 上连续.

对 $i=1,\cdots,n-1$, 在 $B^-$ 上有
\begin{align*}
D_i(\mathcal{R}f)(x)
=-3D_if(R_1x)+4D_if(R_2x).
\end{align*}
对法向导数, 在 $B^-$ 上有
\begin{align*}
D_n(\mathcal{R}f)(x)
=3D_nf(R_1x)-2D_nf(R_2x).
\end{align*}
因此, 当 $x\in B^-$ 趋向 $(a',0)\in\Sigma$ 时,
\begin{align*}
D_i(\mathcal{R}f)(x)
&\longrightarrow
(-3+4)D_if(a',0)=D_if(a',0),
\quad i=1,\cdots,n-1,\\
D_n(\mathcal{R}f)(x)
&\longrightarrow
(3-2)D_nf(a',0)=D_nf(a',0).
\end{align*}

还需确认这些极限确实给出 $\Sigma$ 上的偏导数. 切向偏导数直接由
$(\mathcal{R}f)(x',0)=f(x',0)$ 得到. 对法向偏导数, 当 $t<0$ 且 $t\to0$ 时,
\begin{align*}
f\left(a',-t\right)
&=f(a',0)-tD_nf(a',0)+o(|t|),\\
f\left(a',-\frac{t}{2}\right)
&=f(a',0)-\frac{t}{2}D_nf(a',0)+o(|t|).
\end{align*}
代入\eqref{align:ch05-反射延拓}, 得到
\begin{align*}
(\mathcal{R}f)(a',t)
=f(a',0)+tD_nf(a',0)+o(|t|).
\end{align*}
故从下方计算的法向偏导数也是 $D_nf(a',0)$, 与上方一致. 因此 $\mathcal{R}f$ 的所有一阶偏导数在 $B$ 上存在且连续, 即
\begin{align*}
\mathcal{R}f\in C^{1}(B).
\end{align*}

\noindent$\bm{Step}\bm{2}$.{\heiti (建立反射算子的范数估计)}

先设 $1\leqslant p<\infty$. 由\cref{Basis of Analytics-proposition:经典初等不等式},
\begin{align*}
|a+b|^p\leqslant2^{p-1}(|a|^p+|b|^p).
\end{align*}
因此
\begin{align}
\int_{B^-}|\mathcal{R}f(x)|^p\,\mathrm{d}x
&\leqslant
2^{p-1}\left(
3^p\int_{B^-}|f(R_1x)|^p\,\mathrm{d}x
+4^p\int_{B^-}|f(R_2x)|^p\,\mathrm{d}x
\right).
\label{eq::jiojiojiojioqwio2i3333-1}
\end{align}
由于 $R_1$ 将 $B^-$ 映到 $B^+$, 且 $|\det DR_1|=1$, 换元得到
\begin{align}
\int_{B^-}|f(R_1x)|^p\,\mathrm{d}x
=\int_{B^+}|f(y)|^p\,\mathrm{d}y.
\label{eq::jiojiojiojioqwio2i3333-0}
\end{align}
另一方面, $R_2$ 将 $B^-$ 映到
\begin{align*}
H\triangleq
\{y\in\mathbb{R}^{n}\mid |y'|^2+4y_n^2<\rho^2,\ y_n>0\}
\subseteq B^+,
\end{align*}
且 $|\det DR_2|=\frac{1}{2}$. 因而
\begin{align}
\int_{B^-}|f(R_2x)|^p\,\mathrm{d}x
=2\int_H|f(y)|^p\,\mathrm{d}y
\leqslant2\int_{B^+}|f(y)|^p\,\mathrm{d}y.
\label{eq::jiojiojiojioqwio2i3333-2}
\end{align}
记
\begin{align*}
A_p\triangleq1+2^{p-1}(3^p+2\cdot4^p).
\end{align*}
结合\eqref{eq::jiojiojiojioqwio2i3333-1},
\eqref{eq::jiojiojiojioqwio2i3333-0}及
\eqref{eq::jiojiojiojioqwio2i3333-2}, 并计入上半球的积分, 得到
\begin{align}
\|\mathcal{R}f\|_{L^p(B)}^p
\leqslant A_p\|f\|_{L^p(B^+)}^p.
\label{eq:::HJiuhj3u892HFFHJIOEJO}
\end{align}

对梯度, 在 $B^-$ 上有
\begin{align*}
D(\mathcal{R}f)(x)
=-3Df(R_1x)DR_1+4Df(R_2x)DR_2.
\end{align*}
由于 $\|DR_1\|=1$ 且 $\|DR_2\|\leqslant1$, 故
\begin{align*}
|D(\mathcal{R}f)(x)|
\leqslant3|Df(R_1x)|+4|Df(R_2x)|.
\end{align*}
再次使用上述初等不等式及两次换元, 得到
\begin{align*}
\int_{B^-}|D(\mathcal{R}f)|^p\,\mathrm{d}x
\leqslant
2^{p-1}(3^p+2\cdot4^p)
\int_{B^+}|Df|^p\,\mathrm{d}x.
\end{align*}
加上 $B^+$ 上的积分, 得
\begin{align}
\|D(\mathcal{R}f)\|_{L^p(B)}^p
\leqslant A_p\|Df\|_{L^p(B^+)}^p.
\label{eq:::HJiuhj3u892HFFHJIOEJO-1}
\end{align}
因此, 由\eqref{eq:::HJiuhj3u892HFFHJIOEJO}和
\eqref{eq:::HJiuhj3u892HFFHJIOEJO-1},
\begin{align*}
\|\mathcal{R}f\|_{W^{1,p}(B)}
\leqslant A_p^{\frac{1}{p}}
\|f\|_{W^{1,p}(B^+)}.
\end{align*}

对于光滑函数的 $p=\infty$ 估计, 直接由反射公式及梯度公式得到
\begin{align*}
\|\mathcal{R}f\|_{L^\infty(B)}
&\leqslant7\|f\|_{L^\infty(B^+)},\\
\|D(\mathcal{R}f)\|_{L^\infty(B)}
&\leqslant7\|Df\|_{L^\infty(B^+)},
\end{align*}
从而
\begin{align*}
\|\mathcal{R}f\|_{W^{1,\infty}(B)}
\leqslant7\|f\|_{W^{1,\infty}(B^+)}.
\end{align*}
统一记
\begin{align*}
\kappa_p\triangleq
\begin{cases}
A_p^{\frac{1}{p}},&1\leqslant p<\infty,\\
7,&p=\infty.
\end{cases}
\end{align*}
至此, 已经对 $C^1$ 函数证明
\begin{align*}
\|\mathcal{R}f\|_{W^{1,p}(B)}
\leqslant\kappa_p\|f\|_{W^{1,p}(B^+)},
\quad1\leqslant p\leqslant\infty.
\end{align*}
一般 $W^{1,\infty}$ 函数的反射正则性将在后面单独证明.

\noindent$\bm{Step}\bm{3}$.{\heiti (严格构造边界拉直映射)}

固定 $x_0\in\partial U$. 根据\cref{definition:appC-Ck边界}, 存在正交矩阵 $Q$, 半径 $r>0$ 及 $C^1$ 函数 $\gamma$, 使得在坐标
\begin{align}
z\triangleq Q(x-x_0)
\label{eq:ch05-straighten-coord}
\end{align}
下, 若记 $\widetilde{U}\triangleq Q(U-x_0)$, 则
\begin{align}
\widetilde{U}\cap B(0,r)
=
\{z\in B(0,r)\mid z_n>\gamma(z')\}.
\label{eq:ch05-boundary-def}
\end{align}
由 $0\in\partial\widetilde{U}$ 及 $\gamma$ 的连续性可知 $\gamma(0)=0$: 若 $\gamma(0)<0$, 则 $0\in\widetilde{U}$; 若 $\gamma(0)>0$, 则原点的某个邻域与 $\widetilde{U}$ 不相交. 两者均与 $0\in\partial\widetilde{U}$ 矛盾.

这里区域位于边界图像的上方, 因而 $z_n$ 的正方向指向区域内部. 若额外选取法向坐标, 则应使 $Q$ 的最后一行为 $-\nu(x_0)^T$, 其中 $\nu(x_0)$ 是单位外法向量.

定义
\begin{align*}
\Phi_1(z)
&\triangleq(z',z_n-\gamma(z')),\\
\Psi_1(y)
&\triangleq(y',y_n+\gamma(y')).
\end{align*}
直接代入可得
\begin{align*}
\Phi_1(\Psi_1(y))=y,\quad
\Psi_1(\Phi_1(z))=z.
\end{align*}
故它们互为逆映射. 其导数矩阵为
\begin{align*}
D\Phi_1(z)
&=
\begin{pmatrix}
I_{n-1}&0\\
-D\gamma(z')&1
\end{pmatrix},\\
D\Psi_1(y)
&=
\begin{pmatrix}
I_{n-1}&0\\
D\gamma(y')&1
\end{pmatrix},
\end{align*}
从而
\begin{align*}
\det D\Phi_1(z)=\det D\Psi_1(y)=1.
\end{align*}

下面明确选取拉直后的球, 同时保证其原像紧包含在 $V$ 中. 记
\begin{align*}
d&\triangleq
\operatorname{dist}(x_0,\mathbb{R}^{n}\setminus V)>0,\\
L&\triangleq
\max\limits_{|z'|\leqslant r}|D\gamma(z')|,\\
\rho&\triangleq
\frac{\min\{r,d\}}{4(1+L)}.
\end{align*}
当 $n=1$ 时, 省去带撇的变量并取 $L=0$.

由 $\gamma(0)=0$ 和微分中值估计, 对 $|y|\leqslant2\rho$, 有
\begin{align*}
|\Psi_1(y)|
&\leqslant |y|+|\gamma(y')|\\
&\leqslant(1+L)|y|
\leqslant\frac{\min\{r,d\}}{2}.
\end{align*}
因此可以令
\begin{align*}
B&\triangleq B(0,\rho),\\
\Psi(y)&\triangleq x_0+Q^T\Psi_1(y),\\
W&\triangleq\Psi(B),\\
\Phi(x)&\triangleq\Phi_1(Q(x-x_0)).
\end{align*}
映射 $\Phi:W\to B$ 与 $\Psi:B\to W$ 是互为逆映射的 $C^1$ 微分同胚, 并且
\begin{align*}
\overline{W}
=\Psi(\overline{B})
\subseteq
\overline{B\left(x_0,\frac{\min\{r,d\}}{2}\right)}
\subseteq B(x_0,r)\cap V.
\end{align*}
特别地, $\overline{W}$ 是 $V$ 的紧子集.

由\cref{theorem:向量求导基本性质}, 有
\begin{align}
\Phi(W)=B,\quad\Psi(B)=W,
\label{eq::jifjiojiojiojiomfijsfiojjsifjioejiojmwfiejfiwmfiowijijwfio-1}\\
|\det D\Phi(x)|
=
|\det D\Phi_1(Q(x-x_0))|\,|\det Q|
=1,
\label{eq::jifjiojiojiojiomfijsfiojjsifjioejiojmwfiejfiwmfiowijijwfio-2}\\
|\det D\Psi(y)|
=
|\det Q^T|\,|\det D\Psi_1(y)|
=1.
\label{eq::jifjiojiojiojiomfijsfiojjsifjioejiojmwfiejfiwmfiowijijwfio-3}
\end{align}
这里使用的是行列式的绝对值; 正交矩阵的行列式可以是 $1$ 或 $-1$.

对 $x\in W$, 令 $z=Q(x-x_0)$. 根据\eqref{eq:ch05-boundary-def},
\begin{align*}
x\in U
&\Longleftrightarrow z_n>\gamma(z')\\
&\Longleftrightarrow \Phi_n(x)>0.
\end{align*}
因此
\begin{align*}
\Phi(W\cap U)=B^+,\quad
\Phi(W\cap\partial U)=\Sigma.
\end{align*}
这就是\cref{figure:ch05-拉直边界}所示的边界拉直.

对矩阵采用欧氏算子范数. 由上述矩阵的具体形式,
\begin{align*}
\|D\Phi_1(z)\|\leqslant1+|D\gamma(z')|,\quad
\|D\Psi_1(y)\|\leqslant1+|D\gamma(y')|.
\end{align*}
结合正交矩阵的算子范数为 $1$, 可知
\begin{align*}
C_1&\triangleq
\max\limits_{x\in\overline{W}}\|D\Phi(x)\|<\infty,\\
C_2&\triangleq
\max\limits_{y\in\overline{B}}\|D\Psi(y)\|<\infty,\\
C_3&\triangleq\max\{1,C_1,C_2\}<\infty.
\end{align*}
这些最大值存在, 是因为导数矩阵在相应紧集上连续, 而算子范数是连续函数; 相关算子范数性质可参见
\cref{Functional Analysis-theorem:有界线性算子全体构成的空间完备,Functional Analysis-proposition:有界线性算子的基本不等式}.

\noindent$\bm{Step}\bm{4}$.{\heiti (构造光滑函数的局部扩张并处理坐标变换下的范数)}

先设 $u\in C^1(\overline{U})$. 定义
\begin{align*}
u'(y)\triangleq u(\Psi(y)),\quad y\in B^+,
\end{align*}
并令
\begin{align*}
\overline{u}'\triangleq\mathcal{R}u',\quad
\overline{u}(x)\triangleq\overline{u}'(\Phi(x)),\quad x\in W.
\end{align*}
由前两步,
\begin{align}
\overline{u}'=u'\quad\hbox{在 }B^+\hbox{ 上},
\label{eq:ch05-ext-in-ball-1}\\
\|\overline{u}'\|_{W^{1,p}(B)}
\leqslant\kappa_p\|u'\|_{W^{1,p}(B^+)},
\quad1\leqslant p\leqslant\infty.
\label{eq:ch05-ext-in-ball-2}
\end{align}
此外, $\overline{u}\in C^1(W)$, 且对 $x\in W\cap U$,
\begin{align*}
\overline{u}(x)
=u'(\Phi(x))
=u(\Psi(\Phi(x)))
=u(x).
\end{align*}

先处理 $1\leqslant p<\infty$. 由换元公式及
\eqref{eq::jifjiojiojiojiomfijsfiojjsifjioejiojmwfiejfiwmfiowijijwfio-1}--\eqref{eq::jifjiojiojiojiomfijsfiojjsifjioejiojmwfiejfiwmfiowijijwfio-3},
\begin{align}
\int_W|\overline{u}(x)|^p\,\mathrm{d}x
=
\int_B|\overline{u}'(y)|^p\,\mathrm{d}y,
\label{eq::jijioqwio93283r3r23}\\
\int_{B^+}|u'(y)|^p\,\mathrm{d}y
=
\int_{W\cap U}|u(x)|^p\,\mathrm{d}x.
\label{eq::jijioqwio93283r3r23-111}
\end{align}
链式法则给出
\begin{align}
D\overline{u}(x)
=D\overline{u}'(\Phi(x))D\Phi(x).
\label{eq:ch05-chain-rule-uv}
\end{align}
由算子范数的基本不等式,
\begin{align}
|D\overline{u}(x)|
\leqslant C_1|D\overline{u}'(\Phi(x))|,
\label{eq::jioJFIOWJIOIJOjiojfwio-1}\\
|Du'(y)|
=|Du(\Psi(y))D\Psi(y)|
\leqslant C_2|Du(\Psi(y))|.
\label{eq::jioJFIOWJIOIJOjiojfwio-2}
\end{align}
再次换元, 得到
\begin{align}
\int_W|D\overline{u}(x)|^p\,\mathrm{d}x
\leqslant C_1^p
\int_B|D\overline{u}'(y)|^p\,\mathrm{d}y,
\label{eq::jijioqwio93283r3r23-1}\\
\int_{B^+}|Du'(y)|^p\,\mathrm{d}y
\leqslant C_2^p
\int_{W\cap U}|Du(x)|^p\,\mathrm{d}x.
\label{eq::jijioqwio93283r3r23-2}
\end{align}
将函数及梯度的估计合并, 得到
\begin{align}
\|\overline{u}\|_{W^{1,p}(W)}
\leqslant C_3\|\overline{u}'\|_{W^{1,p}(B)},
\label{eq:ch05-local-est-1}\\
\|u'\|_{W^{1,p}(B^+)}
\leqslant C_3\|u\|_{W^{1,p}(W\cap U)}.
\label{eq:ch05-local-est-2}
\end{align}
因此, 结合\eqref{eq:ch05-ext-in-ball-2},
\begin{align}
\|\overline{u}\|_{W^{1,p}(W)}
\leqslant
\kappa_p C_3^2
\|u\|_{W^{1,p}(W\cap U)}.
\label{eq:ch05-local-est-final}
\end{align}

下面说明 $L^\infty$ 范数在这些坐标变换下的性质. 对任意可测集 $A\subseteq B$, 换元公式给出
\begin{align*}
|\Psi(A)|
=
\int_A|\det D\Psi(y)|\,\mathrm{d}y
=|A|.
\end{align*}
对任意可测集 $F\subseteq W$, 也有
\begin{align*}
|\Phi(F)|
=
\int_F|\det D\Phi(x)|\,\mathrm{d}x
=|F|.
\end{align*}
这里 $|\cdot|$ 表示 Lebesgue 测度. 特别地, $\Phi$ 和 $\Psi$ 都保持零测集, 并且一个集合是零测集当且仅当其对应的像是零测集.

因此, 对任意可测函数 $f:B\to\mathbb{R}$ 及任意 $M\geqslant0$,
$f$ 在 $B$ 上几乎处处满足 $|f|\leqslant M$, 当且仅当
$f\circ\Phi$ 在 $W$ 上几乎处处满足同一不等式. 从本性上确界的定义便得到
\begin{align}
\|f\circ\Phi\|_{L^\infty(W)}
\leqslant\|f\|_{L^\infty(B)},
\label{eq::jiojfwiojiojwio3i2900384}\\
\|f\|_{L^\infty(B)}
\leqslant\|f\circ\Phi\|_{L^\infty(W)}.
\label{eq::jiojfwiojiojwio3i2900382}
\end{align}
这些不等式也包括范数为无穷大的情形. 合并
\eqref{eq::jiojfwiojiojwio3i2900384}和
\eqref{eq::jiojfwiojiojwio3i2900382}, 并将同一论证用于
$\Psi:B^+\to W\cap U$, 可得
\begin{align}
\|f\circ\Phi\|_{L^\infty(W)}
&=\|f\|_{L^\infty(B)},
\nonumber\\
\|g\circ\Psi\|_{L^\infty(B^+)}
&=\|g\|_{L^\infty(W\cap U)}.
\label{eq::jiowiowijoijoijijojsiosioios3829237489}
\end{align}

于是对当前的光滑函数, 由\eqref{eq::jioJFIOWJIOIJOjiojfwio-1},
\eqref{eq::jioJFIOWJIOIJOjiojfwio-2}及
\eqref{eq::jiowiowijoijoijijojsiosioios3829237489},
\begin{align}
\|\overline{u}\|_{W^{1,\infty}(W)}
&\leqslant
\|\overline{u}'\|_{L^\infty(B)}
+C_1\|D\overline{u}'\|_{L^\infty(B)}
\nonumber\\
&\leqslant C_3\|\overline{u}'\|_{W^{1,\infty}(B)},
\label{eq::jiwjiofijoswjiowiow2390jw89023890}\\
\|u'\|_{W^{1,\infty}(B^+)}
&\leqslant
\|u\|_{L^\infty(W\cap U)}
+C_2\|Du\|_{L^\infty(W\cap U)}
\nonumber\\
&\leqslant C_3\|u\|_{W^{1,\infty}(W\cap U)}.
\label{eq::jiwjiofijoswjiowiow2390jw89023890-1}
\end{align}
结合反射估计, 得到与有限 $p$ 情形一致的结论: 对光滑函数 $u$, 其局部扩张满足
\begin{align}
\overline{u}|_{W\cap U}=u,\quad
\|\overline{u}\|_{W^{1,p}(W)}
\leqslant
\kappa_p C_3^2\|u\|_{W^{1,p}(W\cap U)},
\quad1\leqslant p\leqslant\infty.
\label{eq::eEREFFwww}
\end{align}
此时并不要求局部扩张 $\overline{u}$ 在 $W$ 中有紧支集. 紧支集将在乘上光滑截断函数后获得.

\noindent$\bm{Step}\bm{5}$.{\heiti (固定有限覆盖并构造光滑函数的整体扩张)}

由 $\partial U$ 的紧性, 可以从上述边界邻域中选出有限个
$W_1,\cdots,W_N$, 使得
\begin{align*}
\partial U\subseteq\bigcup\limits_{i=1}^{N}W_i,\quad
\overline{W_i}\subseteq V.
\end{align*}
固定各邻域上的坐标映射及反射构造, 将相应局部扩张记为 $\overline{u}_i$. 记
\begin{align*}
c_{i,p}\triangleq\kappa_p C_{3,i}^{\,2},
\quad i=1,\cdots,N.
\end{align*}
由\eqref{eq::eEREFFwww},
\begin{align}
\|\overline{u}_i\|_{W^{1,p}(W_i)}
\leqslant
c_{i,p}\|u\|_{W^{1,p}(W_i\cap U)}
\leqslant
c_{i,p}\|u\|_{W^{1,p}(U)}.
\label{eq::eEREFFwww1}
\end{align}

根据\cref{Real Analysis-theorem:边界和开集基本性质,Real Analysis-theorem:边界和开集基本性质-5},
$\overline{U}=U\cup\partial U$. 由于所有边界点均被上述邻域覆盖, 集合
\begin{align}
K\triangleq
U\setminus\bigcup\limits_{i=1}^{N}W_i
=
\overline{U}\setminus\bigcup\limits_{i=1}^{N}W_i
\label{eq:ioujiwjioKKKI}
\end{align}
是 $U$ 的紧子集.

若 $K\neq\varnothing$, 由
\cref{Real Analysis-proposition:点集间的距离的性质,Real Analysis-proposition:点集间的距离的性质-3}, 可取
\begin{align*}
a\triangleq
\frac{1}{2}\operatorname{dist}(K,\mathbb{R}^{n}\setminus U)>0,\quad
W_0\triangleq
\{x\in\mathbb{R}^{n}\mid\operatorname{dist}(x,K)<a\}.
\end{align*}
于是 $K\subseteq W_0$, 且 $\overline{W_0}$ 是 $U$ 的紧子集.
若 $K=\varnothing$, 则任取一个闭包包含在 $U$ 内的非空开球作为 $W_0$.
两种情况下均有
\begin{align*}
\overline{U}\subseteq\bigcup\limits_{i=0}^{N}W_i.
\end{align*}

根据\cref{lemma:ch05-光滑单位分解-p7R3}, 选取并固定
$\zeta_i\in C_c^\infty(W_i)$, 使得
\begin{align*}
0\leqslant\zeta_i\leqslant1,\quad
\sum\limits_{i=0}^{N}\zeta_i=1
\end{align*}
在包含 $\overline{U}$ 的某个开邻域上成立.
定义
\begin{align*}
K_*\triangleq
\bigcup\limits_{i=0}^{N}\operatorname{spt}\zeta_i.
\end{align*}
这是包含在 $V$ 中的固定紧集.

下面先记录一个适用于所有 $1\leqslant p\leqslant\infty$ 的截断及零延拓结论.
若 $w\in W^{1,p}(W_i)$, 定义
\begin{align*}
g(x)\triangleq
\begin{cases}
\zeta_i(x)w(x),&x\in W_i,\\
0,&x\notin W_i.
\end{cases}
\end{align*}
对任意 $\varphi\in C_c^\infty(\mathbb{R}^{n})$, 有
$\zeta_i\varphi\in C_c^\infty(W_i)$. 因此对 $j=1,\cdots,n$,
\begin{align*}
\int_{\mathbb{R}^{n}}gD_j\varphi\,\mathrm{d}x
&=\int_{W_i}w\zeta_iD_j\varphi\,\mathrm{d}x\\
&=\int_{W_i}wD_j(\zeta_i\varphi)\,\mathrm{d}x
-\int_{W_i}w(D_j\zeta_i)\varphi\,\mathrm{d}x\\
&=-\int_{W_i}
\bigl(\zeta_iD_jw+wD_j\zeta_i\bigr)\varphi\,\mathrm{d}x.
\end{align*}
故 $g$ 的弱导数等于
$\zeta_iDw+wD\zeta_i$ 在 $W_i$ 上的值作零延拓. 这也具体验证了
\cref{theorem:ch05-弱导数性质-4}在本处的应用, 并说明零延拓不会产生额外的边界项.

记
\begin{align*}
a_i&\triangleq\|\zeta_i\|_{L^\infty(W_i)},\\
b_i&\triangleq\|D\zeta_i\|_{L^\infty(W_i)},\\
M_i&\triangleq2a_i+b_i.
\end{align*}
这些常数均有限. 对所有 $1\leqslant p\leqslant\infty$, 有
\begin{align*}
\|g\|_{L^p(\mathbb{R}^{n})}
&\leqslant a_i\|w\|_{L^p(W_i)},\\
\|Dg\|_{L^p(\mathbb{R}^{n})}
&\leqslant
a_i\|Dw\|_{L^p(W_i)}
+b_i\|w\|_{L^p(W_i)}.
\end{align*}
从而
\begin{align*}
\|g\|_{W^{1,p}(\mathbb{R}^{n})}
\leqslant M_i\|w\|_{W^{1,p}(W_i)}.
\end{align*}

现在对 $u\in C^1(\overline{U})$, 令
\begin{align*}
\overline{u}_0\triangleq u|_{W_0},\quad c_{0,p}\triangleq1,
\end{align*}
并将每个 $\zeta_i\overline{u}_i$ 按上述方式作零延拓. 定义
\begin{align*}
\overline{u}\triangleq
\sum\limits_{i=0}^{N}\zeta_i\overline{u}_i.
\end{align*}
由单位分解及局部扩张性质, 对 $x\in U$,
\begin{align*}
\overline{u}(x)
=
\sum\limits_{i=0}^{N}\zeta_i(x)u(x)
=u(x).
\end{align*}
而且
\begin{align*}
\operatorname{spt}\overline{u}
\subseteq K_*\subseteq V.
\end{align*}
令
\begin{align*}
C_{E,p}\triangleq
\sum\limits_{i=0}^{N}M_i c_{i,p}.
\end{align*}
由截断估计和\eqref{eq::eEREFFwww1},
\begin{align}
\|\overline{u}\|_{W^{1,p}(\mathbb{R}^{n})}
&\leqslant
\sum\limits_{i=0}^{N}
M_i\|\overline{u}_i\|_{W^{1,p}(W_i)}
\nonumber\\
&\leqslant C_{E,p}\|u\|_{W^{1,p}(U)}.
\label{eq::fjiwjio233--12}
\end{align}
上述所有变换、反射和乘截断函数的操作都是线性的.

对 $u\in C^\infty(\overline{U})$, 定义 $E_0u\triangleq\overline{u}$. 则 $E_0$ 是线性的, 并且
\begin{align}
\|E_0u\|_{W^{1,p}(\mathbb{R}^{n})}
\leqslant C_{E,p}\|u\|_{W^{1,p}(U)},
\quad1\leqslant p\leqslant\infty.
\label{eqJIOFJIOWJIO323u89rfu3}
\end{align}

\noindent$\bm{Step}\bm{6}$.{\heiti (完成有限 $p$ 情形的扩张)}

设 $1\leqslant p<\infty$, $u\in W^{1,p}(U)$.
由\cref{theorem:ch05-直到边界光滑的函数整体逼近定理}, 可以选取
$u_m\in C^\infty(\overline{U})$, 使得
\begin{align*}
\lim\limits_{m\to\infty}
\|u_m-u\|_{W^{1,p}(U)}=0.
\end{align*}
由 $E_0$ 的线性性及\eqref{eqJIOFJIOWJIO323u89rfu3},
\begin{align*}
\|E_0u_m-E_0u_l\|_{W^{1,p}(\mathbb{R}^{n})}
\leqslant
C_{E,p}\|u_m-u_l\|_{W^{1,p}(U)}.
\end{align*}
故 $\{E_0u_m\}$ 是 $W^{1,p}(\mathbb{R}^{n})$ 中的 Cauchy 列.
由\cref{theorem:ch05-WkpBanach}, 存在该空间中的极限, 定义
\begin{align*}
Eu\triangleq
\lim\limits_{m\to\infty}E_0u_m
\quad\hbox{于 }W^{1,p}(\mathbb{R}^{n}).
\end{align*}

若另一列 $v_m\in C^\infty(\overline{U})$ 也在 $W^{1,p}(U)$ 中趋于 $u$, 则
\begin{align*}
\|E_0u_m-E_0v_m\|_{W^{1,p}(\mathbb{R}^{n})}
&\leqslant
C_{E,p}\|u_m-v_m\|_{W^{1,p}(U)}\\
&\leqslant
C_{E,p}\left(
\|u_m-u\|_{W^{1,p}(U)}
+\|v_m-u\|_{W^{1,p}(U)}
\right)
\longrightarrow0.
\end{align*}
所以两列的扩张具有相同极限, $E$ 的定义与逼近序列无关.

若 $u_m\to u$, $v_m\to v$, 则对任意标量 $a,b$,
$au_m+bv_m\to au+bv$. 利用 $E_0$ 的线性性并取极限, 得
\begin{align*}
E(au+bv)=aEu+bEv.
\end{align*}
再由范数连续性,
\begin{align*}
\|Eu\|_{W^{1,p}(\mathbb{R}^{n})}
&=
\lim\limits_{m\to\infty}
\|E_0u_m\|_{W^{1,p}(\mathbb{R}^{n})}\\
&\leqslant
C_{E,p}
\lim\limits_{m\to\infty}\|u_m\|_{W^{1,p}(U)}
=C_{E,p}\|u\|_{W^{1,p}(U)}.
\end{align*}

由于 $E_0u_m=u_m$ 在 $U$ 上成立,
\begin{align*}
\|Eu-u\|_{L^p(U)}
&\leqslant
\|Eu-E_0u_m\|_{L^p(\mathbb{R}^{n})}
+\|u_m-u\|_{L^p(U)}
\longrightarrow0.
\end{align*}
因此 $Eu=u$ 在 $U$ 上几乎处处成立.

所有 $E_0u_m$ 的支集均包含在同一个紧集 $K_*$ 中, 所以
\begin{align*}
\|Eu\|_{L^p(\mathbb{R}^{n}\setminus K_*)}
\leqslant
\|Eu-E_0u_m\|_{L^p(\mathbb{R}^{n})}
\longrightarrow0.
\end{align*}
故 $Eu=0$ 在 $\mathbb{R}^{n}\setminus K_*$ 上几乎处处成立, 即
\begin{align*}
\operatorname{spt}(Eu)\subseteq K_*\subseteq V.
\end{align*}
有限 $p$ 的情形至此得证.

\noindent$\bm{Step}\bm{7}$.{\heiti (对一般 $W^{1,\infty}$ 函数拉直边界并证明半球上的 Lipschitz 估计)}

现在设 $u\in W^{1,\infty}(U)$.
本情形直接构造扩张, 不使用 $W^{1,\infty}$ 范数下的光滑逼近.

由\cref{theorem:ch05-W1infty局部刻画}, $u$ 有一个局部 Lipschitz 连续的代表元, 以下仍记为 $u$.
这个连续代表元是唯一的: 两个连续代表元若几乎处处相等而在某点不等, 则由连续性可知它们在该点的某个开邻域内仍不等, 与几乎处处相等矛盾.

固定前面选定的一个边界坐标邻域 $W$, 仍记相应映射为
$\Phi:W\to B$ 和 $\Psi:B\to W$. 定义
\begin{align*}
u'(y)\triangleq u(\Psi(y)),\quad y\in B^+.
\end{align*}
由于 $u$ 局部 Lipschitz, $\Psi$ 为 $C^1$ 映射, $u'$ 在 $B^+$ 中局部 Lipschitz.

下面验证其弱导数的具体公式. 局部 Lipschitz 函数几乎处处可微, 且其经典导数与弱导数几乎处处一致. 取零测集 $N\subseteq W\cap U$, 使得 $u$ 在
$(W\cap U)\setminus N$ 上可微, 且经典导数等于弱导数.
由于 $\Phi$ 保持零测集, $\Phi(N)$ 也是零测集.
对 $y\in B^+\setminus\Phi(N)$, 经典链式法则给出
\begin{align*}
Du'(y)=Du(\Psi(y))D\Psi(y).
\end{align*}
又因为 $u'$ 局部 Lipschitz, 此式同时给出其弱导数, 并在 $B^+$ 上几乎处处成立.

结合\eqref{eq::jiowiowijoijoijijojsiosioios3829237489}, 得到
\begin{align*}
\|u'\|_{L^\infty(B^+)}
&=\|u\|_{L^\infty(W\cap U)},\\
\|Du'\|_{L^\infty(B^+)}
&\leqslant
C_2\|Du\|_{L^\infty(W\cap U)}.
\end{align*}
因此 $u'\in W^{1,\infty}(B^+)$, 且
\begin{align}
\|u'\|_{W^{1,\infty}(B^+)}
\leqslant
C_3\|u\|_{W^{1,\infty}(W\cap U)}.
\label{eq:u'-W1infty估计}
\end{align}

为了将反射函数跨越 $\Sigma$ 拼接, 还需要 $u'$ 在整个 $B^+$ 上的 Lipschitz 估计.
这里 $B^+$ 的凸性至关重要. 记
\begin{align*}
H\triangleq\|Du'\|_{L^\infty(B^+)}.
\end{align*}
任取 $y,z\in B^+$, 线段
\begin{align*}
S\triangleq
\{(1-t)y+tz\mid0\leqslant t\leqslant1\}
\end{align*}
是 $B^+$ 的紧子集. 因而
\begin{align*}
d_S\triangleq
\operatorname{dist}(S,\mathbb{R}^{n}\setminus B^+)>0.
\end{align*}
取非负标准磨光核 $\eta$, 满足
\begin{align*}
\eta\in C_c^\infty(B(0,1)),\quad
\int_{\mathbb{R}^{n}}\eta(h)\,\mathrm{d}h=1,\quad
\eta_\varepsilon(h)
=\frac{1}{\varepsilon^n}\eta\left(\frac{h}{\varepsilon}\right).
\end{align*}
当 $0<\varepsilon<\frac{d_S}{2}$ 时, 可以在线段 $S$ 的邻域内定义
\begin{align*}
u'_\varepsilon(\xi)
\triangleq
\int_{\mathbb{R}^{n}}
\eta_\varepsilon(h)u'(\xi-h)\,\mathrm{d}h.
\end{align*}
此时积分只使用 $B^+$ 内部的值, 且由弱导数与内部磨光的关系,
\begin{align*}
Du'_\varepsilon(\xi)
=
\int_{\mathbb{R}^{n}}
\eta_\varepsilon(h)Du'(\xi-h)\,\mathrm{d}h.
\end{align*}
因此在 $S$ 上有
\begin{align*}
|Du'_\varepsilon(\xi)|\leqslant H.
\end{align*}
沿线段应用微积分基本定理, 得
\begin{align*}
|u'_\varepsilon(y)-u'_\varepsilon(z)|
&\leqslant
\int_0^1
|Du'_\varepsilon((1-t)y+tz)|\,|y-z|\,\mathrm{d}t\\
&\leqslant H|y-z|.
\end{align*}
由于 $u'$ 连续, 内部磨光在紧集 $S$ 上一致趋于 $u'$.
令 $\varepsilon\downarrow0$, 得到
\begin{align*}
|u'(y)-u'(z)|\leqslant H|y-z|,
\quad y,z\in B^+.
\end{align*}
所以 $u'$ 在整个 $B^+$ 上是 Lipschitz 函数, 且其 Lipschitz 常数不超过
$\|Du'\|_{L^\infty(B^+)}$.

由此, $u'$ 唯一连续延拓到 $\overline{B^+}$, 并保持同一个 Lipschitz 常数.
同时, 连续性保证
\begin{align*}
|u'(y)|\leqslant\|u'\|_{L^\infty(B^+)},
\quad y\in\overline{B^+}.
\end{align*}

\noindent$\bm{Step}\bm{8}$.{\heiti (完整证明一般 $W^{1,\infty}$ 函数的反射和局部扩张)}

对上一步得到的 Lipschitz 函数 $u'$, 仍按\eqref{align:ch05-反射延拓}定义
\begin{align*}
\overline{u}'(y)\triangleq
\begin{cases}
u'(y),&y\in B^+\cup\Sigma,\\
-3u'(R_1y)+4u'(R_2y),&y\in B^-.
\end{cases}
\end{align*}
由于 $u'$ 已连续延拓到 $\overline{B^+}$, 此定义在 $\Sigma$ 上有意义.
当 $y\in B^-$ 趋向 $a\in\Sigma$ 时,
\begin{align*}
\overline{u}'(y)
\longrightarrow
-3u'(a)+4u'(a)=u'(a),
\end{align*}
故两侧在 $\Sigma$ 上连续衔接.

对 $y,z\in B^+$, 有
\begin{align*}
|\overline{u}'(y)-\overline{u}'(z)|
\leqslant H|y-z|.
\end{align*}
对 $y,z\in B^-$, 利用 $R_1,R_2$ 的 Lipschitz 常数均不超过 $1$, 得到
\begin{align*}
|\overline{u}'(y)-\overline{u}'(z)|
&\leqslant
3|u'(R_1y)-u'(R_1z)|
+4|u'(R_2y)-u'(R_2z)|\\
&\leqslant
3H|R_1y-R_1z|+4H|R_2y-R_2z|\\
&\leqslant7H|y-z|.
\end{align*}
由连续性, 这两个估计分别延续到相应半球与 $\Sigma$ 的并集.

最后, 设 $y\in B^+$, $z\in B^-$.
由于 $B$ 凸, 线段 $[y,z]$ 包含在 $B$ 内, 并与 $\Sigma$ 相交于某点 $a$.
此时
\begin{align*}
|y-a|+|a-z|=|y-z|.
\end{align*}
利用前面两侧的估计, 得
\begin{align*}
|\overline{u}'(y)-\overline{u}'(z)|
&\leqslant
|\overline{u}'(y)-\overline{u}'(a)|
+|\overline{u}'(a)-\overline{u}'(z)|\\
&\leqslant H|y-a|+7H|a-z|\\
&\leqslant7H|y-z|.
\end{align*}
这就证明了 $\overline{u}'$ 在整个 $B$ 上是 Lipschitz 函数, Lipschitz 常数不超过 $7H$.

根据\cref{theorem:ch05-W1infty局部刻画}, Lipschitz 函数的弱导数与其几乎处处存在的经典导数一致. 在可微点处, 导数的算子范数不超过 Lipschitz 常数, 因而
\begin{align*}
\|D\overline{u}'\|_{L^\infty(B)}
\leqslant7H
=7\|Du'\|_{L^\infty(B^+)}.
\end{align*}
反射公式还给出
\begin{align*}
\|\overline{u}'\|_{L^\infty(B)}
\leqslant7\|u'\|_{L^\infty(B^+)}.
\end{align*}
因此
\begin{align}
\overline{u}'|_{B^+}=u',\quad
\overline{u}'\in W^{1,\infty}(B),\quad
\|\overline{u}'\|_{W^{1,\infty}(B)}
\leqslant7\|u'\|_{W^{1,\infty}(B^+)}.
\label{eq:延拓u'-W1infty}
\end{align}
这一论证同时排除了反射平面上出现跳跃或额外分布导数的可能, 因为已经证明了整个球上的 Lipschitz 正则性.

定义
\begin{align*}
\overline{u}(x)\triangleq\overline{u}'(\Phi(x)),
\quad x\in W.
\end{align*}
由于 $\overline{u}'$ Lipschitz 而 $\Phi$ 为 $C^1$ 映射, $\overline{u}$ 在 $W$ 中局部 Lipschitz.
与上一步验证 $u'$ 的弱导数时一样, 利用 $\Psi$ 保持零测集, 经典链式法则在几乎处处的点适用, 得到
\begin{align*}
D\overline{u}(x)
=D\overline{u}'(\Phi(x))D\Phi(x)
\end{align*}
在 $W$ 上几乎处处成立.
因此
\begin{align*}
\|\overline{u}\|_{L^\infty(W)}
&=\|\overline{u}'\|_{L^\infty(B)},\\
\|D\overline{u}\|_{L^\infty(W)}
&\leqslant
C_1\|D\overline{u}'\|_{L^\infty(B)}.
\end{align*}
故 $\overline{u}\in W^{1,\infty}(W)$, 并且由
\eqref{eq:u'-W1infty估计}及\eqref{eq:延拓u'-W1infty},
\begin{align}
\|\overline{u}\|_{W^{1,\infty}(W)}
&\leqslant
C_3\|\overline{u}'\|_{W^{1,\infty}(B)}
\nonumber\\
&\leqslant
7C_3\|u'\|_{W^{1,\infty}(B^+)}
\nonumber\\
&\leqslant
7C_3^2\|u\|_{W^{1,\infty}(W\cap U)}.
\label{eq:局部延拓-W1infty}
\end{align}
同时, 对 $x\in W\cap U$,
\begin{align*}
\overline{u}(x)
=\overline{u}'(\Phi(x))
=u'(\Phi(x))
=u(x).
\end{align*}
因此已经对一般 $W^{1,\infty}$ 函数完成局部扩张, 并证明了所需的局部范数估计.

\noindent$\bm{Step}\bm{9}$.{\heiti (完成 $p=\infty$ 情形的整体扩张并验证线性性)}

使用第五步已经固定的邻域
$W_0,\cdots,W_N$ 及单位分解
$\zeta_0,\cdots,\zeta_N$.
对 $i=1,\cdots,N$, 按第七步和第八步构造局部扩张 $\overline{u}_i$, 并令
\begin{align*}
\overline{u}_0\triangleq u|_{W_0}.
\end{align*}
由\eqref{eq:局部延拓-W1infty},
\begin{align*}
\|\overline{u}_i\|_{W^{1,\infty}(W_i)}
&\leqslant
c_{i,\infty}\|u\|_{W^{1,\infty}(U)},
\quad i=1,\cdots,N,\\
c_{i,\infty}&\triangleq7C_{3,i}^{\,2},\quad
c_{0,\infty}\triangleq1.
\end{align*}
将每个 $\zeta_i\overline{u}_i$ 作零延拓, 记为 $g_i$, 并定义
\begin{align*}
Eu\triangleq\sum\limits_{i=0}^{N}g_i.
\end{align*}

第五步已经通过测试函数计算证明, 对这样的零延拓有
\begin{align*}
Dg_i
=
\zeta_iD\overline{u}_i
+\overline{u}_iD\zeta_i
\end{align*}
在 $W_i$ 上成立, 在 $W_i$ 外为零. 因而
\begin{align*}
\|g_i\|_{L^\infty(\mathbb{R}^{n})}
&\leqslant
a_i\|\overline{u}_i\|_{L^\infty(W_i)},\\
\|Dg_i\|_{L^\infty(\mathbb{R}^{n})}
&\leqslant
a_i\|D\overline{u}_i\|_{L^\infty(W_i)}
+b_i\|\overline{u}_i\|_{L^\infty(W_i)}.
\end{align*}
这说明 $g_i\in W^{1,\infty}(\mathbb{R}^{n})$, 且
\begin{align*}
\|g_i\|_{W^{1,\infty}(\mathbb{R}^{n})}
\leqslant
M_i\|\overline{u}_i\|_{W^{1,\infty}(W_i)}.
\end{align*}
于是
\begin{align}
\|Eu\|_{W^{1,\infty}(\mathbb{R}^{n})}
&\leqslant
\sum\limits_{i=0}^{N}
\|g_i\|_{W^{1,\infty}(\mathbb{R}^{n})}
\nonumber\\
&\leqslant
\sum\limits_{i=0}^{N}
M_i\|\overline{u}_i\|_{W^{1,\infty}(W_i)}
\nonumber\\
&\leqslant
\left(
\sum\limits_{i=0}^{N}M_i c_{i,\infty}
\right)
\|u\|_{W^{1,\infty}(U)}.
\label{eq:整体延拓-W1infty}
\end{align}

对 $x\in U$, 每个满足 $\zeta_i(x)\neq0$ 的项都有
$\overline{u}_i(x)=u(x)$. 因此
\begin{align*}
Eu(x)
=
\sum\limits_{i=0}^{N}\zeta_i(x)u(x)
=u(x).
\end{align*}
并且
\begin{align*}
\operatorname{spt}(Eu)
\subseteq
\bigcup\limits_{i=0}^{N}\operatorname{spt}\zeta_i
=K_*\subseteq V.
\end{align*}

最后验证这一构造在 Sobolev 等价类上良定义且线性.
局部 Lipschitz 代表元由其几乎处处等价类唯一确定; 拉直后的 Lipschitz 函数在
$\overline{B^+}$ 上的连续延拓也由其在 $B^+$ 上的值唯一确定.
因此整个构造与最初代表元的选取无关.
若 $u,v\in W^{1,\infty}(U)$, 则 $au+bv$ 的连续代表元就是相应代表元的线性组合, 而连续延拓、反射、坐标复合、乘截断函数及有限求和都保持线性. 故
\begin{align*}
E(au+bv)=aEu+bEv.
\end{align*}
结合\eqref{eq:整体延拓-W1infty}, $E$ 是有界线性算子.

所有坐标邻域、反射公式及单位分解均只根据 $U,V$ 预先选定, 所以所得常数只依赖于 $n,p,U,V$, 与 $u$ 无关.
定理对 $1\leqslant p\leqslant\infty$ 全部成立.
\end{proof}



\end{document}